\begin{summarybox}
\textbf{C010\quad 一页快记：椭圆切线距离}

\textbf{【前提】} $\frac{x^2}{a^2}+\frac{y^2}{b^2}=1$，$a>b>0$，$c^2=a^2-b^2$；$l$ 必须是\textbf{切线}。

\textbf{【三条核心公式】}
\[
\boxed{d_1d_2=b^2},\qquad
\boxed{d_1+d_2=2d_0},\qquad
\boxed{OH_1=OH_2=a}.
\]
垂足完整轨迹：$\boxed{x^2+y^2=a^2}$。中心距范围：$\boxed{b\le d_0\le a}$。

\textbf{【最快推导：3 个信号】}
\[
\begin{gathered}
l:ux+vy=h,\quad u^2+v^2=1,
\quad h^2=a^2u^2+b^2v^2,\\
d_1=h+cu,\quad d_2=h-cu,\\
\boxed{d_1d_2=h^2-c^2u^2=b^2}.
\end{gathered}
\]
\textbf{【最快几何记忆】}焦点关于切线反射 $\to$ 反射点到另一焦点长 $2a$ $\to$ 中位线长 $a$（主圆）；两条中心对称的切线 $\to$ 主圆中同一焦点的一条弦 $\to$ 点幂 $a^2-c^2=b^2$。

\textbf{【考场选择方法】}
\begin{itemize}
\item \textbf{只问乘积}：直接写 $b^2$，不设切线。
\item \textbf{问两个具体距离}：联用 $d_1+d_2=2d_0$、$d_1d_2=b^2$，解一元二次方程。
\item \textbf{问垂足轨迹或构造切线}：先换成主圆 $x^2+y^2=a^2$，必要时做逆向构造。
\item \textbf{已知斜率}：$m^2=a^2k^2+b^2$；注意竖直切线单独处理。
\item \textbf{需要一法通吃}：用 $Ax+By=D$，切线条件 $D^2=a^2A^2+b^2B^2$。
\end{itemize}

\textbf{【记忆口诀】}\quad\textbf{积看短轴，垂足长圆；和为两倍中心距，先看是否真相切。}
\end{summarybox}
